% ───────── Bìa chính ─────────
\begin{titlepage}
\begin{tikzpicture}[remember picture,overlay]
  \draw[line width=2.2pt] ($(current page.north west)+(2.7cm,-1.9cm)$) rectangle ($(current page.south east)+(-1.3cm,1.9cm)$);
\end{tikzpicture}
\begin{center}
{\large ĐẠI HỌC QUỐC GIA HÀ NỘI}\\
{\large\bfseries TRƯỜNG ĐẠI HỌC KHOA HỌC TỰ NHIÊN}\\
\rule{3cm}{0.6pt}

\vspace{4.5cm}
{\bfseries \tacgia}

\vspace{2.8cm}
\begin{minipage}{0.95\textwidth}\centering\large\bfseries\setstretch{1.4}
\tendetai
\end{minipage}

\vspace{2.8cm}
{\bfseries LUẬN VĂN THẠC SĨ KHOA HỌC}

\vfill
{\bfseries Hà Nội - \nam}
\end{center}
\end{titlepage}

% ───────── Bìa phụ ─────────
\begin{titlepage}
\begin{tikzpicture}[remember picture,overlay]
  \draw[line width=2.2pt] ($(current page.north west)+(2.7cm,-1.9cm)$) rectangle ($(current page.south east)+(-1.3cm,1.9cm)$);
\end{tikzpicture}
\begin{center}
{\large ĐẠI HỌC QUỐC GIA HÀ NỘI}\\
{\large\bfseries TRƯỜNG ĐẠI HỌC KHOA HỌC TỰ NHIÊN}\\
\rule{3cm}{0.6pt}

\vspace{3cm}
{\bfseries \tacgia}

\vspace{1.6cm}
\begin{minipage}{0.95\textwidth}\centering\large\bfseries\setstretch{1.4}
\tendetai
\end{minipage}
\end{center}

\vspace{1.2cm}
\hspace{2cm}Chuyên ngành: Khoa học dữ liệu

\vspace{0.2cm}
\hspace{2cm}Mã số: 8460108

\begin{center}
\vspace{1.4cm}
{\bfseries LUẬN VĂN THẠC SĨ KHOA HỌC}
\end{center}

\vspace{1.2cm}
\begin{flushright}
\begin{tabular}{r}
NGƯỜI HƯỚNG DẪN KHOA HỌC:\\[10pt]
\huongdanmot\\[2pt]
\huongdanhai
\end{tabular}
\end{flushright}
\vfill
\begin{center}{\bfseries Hà Nội - \nam}\end{center}
\end{titlepage}

% ───────── Lời cam đoan ─────────
\chapter*{LỜI CAM ĐOAN}
Tôi xin cam đoan luận văn \textbf{\tendetaithuong} là công trình nghiên cứu do tôi thực hiện dưới sự hướng dẫn khoa học của \huongdanmot{} và \huongdanhai, trên cơ sở vận dụng các kiến thức lý thuyết, phương pháp nghiên cứu và tài liệu chuyên ngành có liên quan. Tôi xin cam đoan các nội dung, số liệu và kết quả nghiên cứu được trình bày trong luận văn là trung thực, được thu thập, xử lý và phân tích một cách nghiêm túc. Các tài liệu, công trình nghiên cứu và kết quả của các tác giả khác được sử dụng trong luận văn đều được trích dẫn và ghi nguồn đầy đủ theo đúng quy định.

Một phần nội dung của luận văn đã được công bố trong bài báo khoa học mà tôi là đồng tác giả \cite{tran2026dynamic}; các nội dung kế thừa và các điểm khác biệt so với bài báo đều được chỉ rõ trong luận văn. Tôi xin cam đoan luận văn này chưa từng được sử dụng để bảo vệ nhằm nhận học vị tại bất kỳ cơ sở đào tạo nào khác. Nếu phát hiện có bất kỳ sự sai lệch hoặc gian lận nào trong nội dung luận văn, tôi xin hoàn toàn chịu trách nhiệm trước nhà trường và pháp luật.

\vspace{1cm}
\begin{flushright}
\begin{minipage}{7cm}\centering
\textit{Hà Nội, ngày ... tháng ... năm \nam}\\
\textbf{Học viên}\\[2.6cm]
\textbf{\tacgiathuong}
\end{minipage}
\end{flushright}

% ───────── Lời cảm ơn ─────────
\chapter*{LỜI CẢM ƠN}
Trước hết, tôi xin bày tỏ lòng biết ơn sâu sắc tới \huongdanmot{} và \huongdanhai, những người đã trực tiếp hướng dẫn, định hướng học thuật và đưa ra nhiều góp ý quý báu trong suốt quá trình tôi thực hiện nghiên cứu và hoàn thiện luận văn này. Sự tận tình chỉ bảo của các Thầy Cô đã giúp tôi hình thành tư duy nghiên cứu, khắc phục nhiều hạn chế về phương pháp và trình bày kết quả một cách khoa học, trung thực hơn.

Tôi xin trân trọng cảm ơn các thầy, cô Khoa Toán - Cơ - Tin học, Trường Đại học Khoa học Tự nhiên, Đại học Quốc gia Hà Nội đã giảng dạy, hỗ trợ và tạo điều kiện thuận lợi cho tôi trong quá trình học tập và nghiên cứu. Tôi xin cảm ơn các thầy, cô phản biện đã dành thời gian đọc, nhận xét và đóng góp những ý kiến quý báu, giúp tôi tiếp tục hoàn thiện nội dung luận văn.

Cuối cùng, tôi xin gửi lời cảm ơn tới gia đình, bạn bè và đồng nghiệp đã luôn động viên, chia sẻ và hỗ trợ tôi trong suốt quá trình học tập, nghiên cứu và hoàn thành luận văn. Tôi xin chân thành cảm ơn!

\vspace{1cm}
\begin{flushright}
\begin{minipage}{7cm}\centering
\textit{Hà Nội, ngày ... tháng ... năm \nam}\\
\textbf{Học viên}\\[2.6cm]
\textbf{\tacgiathuong}
\end{minipage}
\end{flushright}

% ───────── Tóm tắt ─────────
\chapter*{TÓM TẮT}
{\small\setstretch{1.15}
\noindent\textbf{Tóm tắt:} Tối ưu hóa danh mục đầu tư là bài toán khó do quan hệ phụ thuộc giữa các tài sản thay đổi theo thời gian và điều kiện thị trường bất định. Các phương pháp truyền thống như mô hình trung bình--phương sai của Markowitz dựa trên ước lượng tĩnh của ma trận hiệp phương sai, còn các mô hình học sâu chuỗi thời gian thường xử lý từng cổ phiếu một cách độc lập, bỏ qua cấu trúc quan hệ giữa các cổ phiếu. Luận văn này nghiên cứu tính khả thi của một mô hình học sâu lai, huấn luyện đầu--cuối, khai thác quan hệ tương quan dài hạn giữa các cổ phiếu để phân bổ danh mục cổ phiếu vốn hóa vừa và lớn tại thị trường chứng khoán Việt Nam.

Mô hình đề xuất gồm ba thành phần: (i) bộ mã hóa thời gian kiểu Transformer với mặt nạ nhân quả, biến chuỗi 60 phiên của 19 đặc trưng (giá và kỹ thuật, tài chính doanh nghiệp, vĩ mô) thành véc-tơ biểu diễn cho từng cổ phiếu; (ii) mạng chú ý đồ thị (Graph Attention Network -- GAT) lan truyền thông tin trên đồ thị tương quan lợi suất dài hạn, ước lượng trên cửa sổ 250 phiên; và (iii) tác tử học tăng cường Proximal Policy Optimization (PPO) với chính sách Dirichlet, được thưởng bằng tỷ số Sharpe vi phân (Differential Sharpe Ratio) tính trên lợi suất ròng sau phí giao dịch. Dữ liệu gồm 222 mã cổ phiếu giai đoạn 2016--2025, chia theo thời gian thành tập huấn luyện (2016--2022), kiểm định (2023) và kiểm tra (2024--2025).

Trước khi xây dựng mô hình, luận văn kiểm chứng rằng quan hệ tương quan dài hạn tồn tại và bền vững: các nhóm cổ phiếu tìm được từ tương quan lợi suất 2016--2022 trùng với các nhóm ngành, 76\% cặp tương quan mạnh tiếp tục tương quan mạnh ở 2024--2025, và đồ thị 250 phiên giữ được 81\% tập cạnh qua mỗi tháng so với 52\% của đồ thị 60 phiên. Trên giai đoạn kiểm tra 2024--2025, mô hình với đồ thị dài hạn đạt tỷ số Sharpe $1{,}25\pm0{,}19$, lợi suất tích lũy 42,2\% và mức sụt giảm tối đa $-7{,}7\%$ trên 10 hạt giống ngẫu nhiên, tốt hơn danh mục phân bổ đều (0,90; 28,0\%; $-8{,}9\%$) và trung vị danh mục ngẫu nhiên ở mọi hạt giống, nhưng chưa vượt VN-Index (1,45) và danh mục trung bình--phương sai (1,56). Trên kiểm định cuốn chiếu 71 tháng (2020--2025, huấn luyện lại mỗi năm), mô hình đạt tỷ số Sharpe $0{,}63\pm0{,}19$ trên 8 hạt giống, cao hơn một chút so với phân bổ đều (0,55), VN-Index (0,54) và danh mục trung bình--phương sai (0,49), dù các chênh lệch chưa có ý nghĩa thống kê. Đồ thị dài hạn là yếu tố có ảnh hưởng lớn và nhất quán nhất: nó cho kết quả tốt hơn đồ thị ngắn hạn trên cả hai giai đoạn và trên tập cổ phiếu thanh khoản tốt, với mức luân chuyển thấp hơn, thông qua việc ưu tiên các cổ phiếu ít liên kết với thị trường. Đây là kết luận duy nhất có ý nghĩa thống kê theo bootstrap khối dừng ($p = 0{,}002$ và $0{,}01$) và không phụ thuộc vào một lựa chọn ngưỡng cụ thể. Học tăng cường tốt hơn học có giám sát trên cùng kiến trúc và hơn một quy tắc không học giữ các cổ phiếu ít liên kết nhất trên đồ thị dài hạn, nhưng chênh lệch so với quy tắc chưa có ý nghĩa thống kê, và trên tập cổ phiếu thanh khoản tốt quy tắc vượt mô hình rõ rệt (1,38 so với 0,80); hai sửa đổi phần thưởng -- học với phí phụ thuộc thanh khoản và thưởng trên lợi suất vượt quy tắc -- không cải thiện được điều này một cách nhất quán. Khoảng cách với VN-Index trên 2024--2025 phần lớn do họ Vingroup: khi loại nhóm này, chỉ số vốn hóa tương ứng có tỷ số Sharpe 0,95, thấp hơn mô hình. Khi giới hạn ở cổ phiếu thanh khoản tốt hoặc khi phí giao dịch từ 30 điểm cơ bản, ưu thế của mô hình thu hẹp đáng kể. Luận văn kết luận rằng phương pháp khả thi và có cơ sở thực nghiệm, đồng thời chỉ rõ các hạn chế về thanh khoản, ý nghĩa thống kê và các hướng cải tiến.

\medskip
\noindent\textbf{\textit{Từ khóa:}} \textit{Tối ưu hóa danh mục đầu tư, tương quan dài hạn, mạng nơ-ron đồ thị, Graph Attention Network, học tăng cường, Proximal Policy Optimization, tỷ số Sharpe vi phân, thị trường chứng khoán Việt Nam.}
}
